\section{Frequency regimes, skin depth, and conductor-mode adaptivity}
\label{sec:frequency-regimes}

Wireless-power systems often operate in a regime where the device is
electrically small while conductor skin/proximity effects are not negligible.
The vNext formulation separates these two notions.

\subsection{Electromagnetic size}

Define a scene electrical-size indicator
\[
 \kappa_{\rm scene}
 =
 \max_{\mathbf r,\mathbf r'\in\mathcal S}
 |k_b|\,\norm{\mathbf r-\mathbf r'} ,
\]
where $k_b$ is the relevant background wavenumber. When
$\kappa_{\rm scene}\ll1$, a magnetoquasistatic or electroquasistatic expansion
may provide an efficient analytic baseline. Correctness of the REFERENCE model,
however, does not rely on dropping retardation unless the artifact explicitly
declares an MQS-only domain.

\subsection{Skin depth}

For a good isotropic conductor, the characteristic skin depth is
\[
 \delta
 \approx
 \sqrt{\frac{2}{\omega\mu\sigma}}.
\]
Define dimensionless cross-section indicators
\[
 \chi_w=\frac{w}{\delta},\qquad
 \chi_h=\frac{h}{\delta}.
\]
These, together with separation-to-size ratios, drive the required
cross-section modal rank.

\subsection{Adaptive modal family}

The conductor basis contains nested mode sets
\[
 V_\perp^{(0)}
 \subset V_\perp^{(1)}
 \subset\cdots .
\]
A teacher solve begins with a rank predicted from
$(\chi_w,\chi_h)$ and geometric proximity, then enriches until
\[
 \eta_{\rm mode}
 =
 \frac{\norm{Z^{(r+\Delta r)}-Z^{(r)}}_{W_Z}}
      {\max(\norm{Z^{(r+\Delta r)}}_{W_Z},\epsilon_Z)}
 \le\tau_{\rm mode}
\]
and selected local current/loss observables also converge.

High-conductivity accuracy is therefore purchased by additional intrinsic
cross-section modes, not by shrinking cells throughout surrounding space.

\subsection{DC and low-frequency limit}

As $\omega\to0$,
\[
 \delta\to\infty,
\]
and the current basis naturally collapses toward the uniform conduction mode.
Charge/current scaling and the mixed formulation must remain well conditioned
in this limit. A DC resistance benchmark is part of the mandatory teacher test
suite.

\subsection{Full-wave transition}

As $\kappa_{\rm scene}$ approaches unity, retardation and radiation are retained
by the full Green kernel. The same continuous scene and port definitions are
used; only kernel evaluation, basis requirements, and acceleration strategy
change.

\subsection{Analytic baseline selection}

The physics baseline used by the neural model is selected by dimensionless
regime indicators:
\[
 \eta_{\rm regime}
 =
 (\kappa_{\rm scene},\chi_w,\chi_h,
  d_{\min}/L_0,\text{material contrast}).
\]
Examples:
\begin{itemize}
\item filament/mutual-inductance baseline for very thin, weakly coupled
      conductors;
\item finite-cross-section MQS/PEEC baseline for electrically small WPT;
\item retarded integral baseline for larger electrical size;
\item layered Green baseline near planar stratified media.
\end{itemize}
The neural correction always receives the baseline identity/quality features so
it does not have to infer which asymptotic model was used.

\subsection{No regime discontinuity}

When multiple baselines are blended, the blend uses smooth regime weights and
the final structured decoder. Switching a backend must not introduce a
discontinuity in $Z$, loss, or design gradients as geometry/frequency crosses a
threshold.
